% 七、预期成果和考核指标。
% 内容迁自 NSFC_2026_MineChangeTrust；保留原稿事实、暂定值和待补项。
\subsection{预期研究结果}
预期形成一致性约束适用条件与误报—召回权衡的实证结论，区分区域选择、强度和判决工作点的解释，并说明残留错误可被复核发现的范围。交付一套可复现协议与结果包，包含四组对照及DonaNet结果、强度和阈值敏感性分析、固定面积复核曲线及失败案例；逐项列明有效矿区数、效应区间和未完成条件，不以预期涨幅代替考核。

按指南要求发表至少1篇SCI/EI源刊论文，提交中期进展及结题报告；投稿、录用以及EI会议论文均不能代替源刊论文发表要求。同时整理可复现的实验与评测流程，列出案例来源并说明各案例能够支持的判断。公开代码和衍生产物前，核对数据许可、既有项目权属及实验室成果共享约定。矿企现场部署、环境风险分级系统和新疆示范区建设不列为本项目的必达成果。
